\section{Introduction}
\label{s-introduction}

Let \(\sigma(n)\) be the sum of the positive divisors of \(n\), and let
\(\gamma\) be Euler's constant. Gr\"onwall \cite{gronwall1913} proved
\[
    \limsup_{n\to\infty}
    \frac{\sigma(n)}{e^\gamma n\log\log n}
    =
    1.
\]
Robin \cite{robin1984} sharpened this maximal-order statement to a
pointwise criterion. The Riemann hypothesis holds if and only if
\[
    \sigma(n)<e^\gamma n\log\log n
    \qquad (n\geq5041).
\]
We maximize the normalized ratio over integers with a fixed number of
distinct prime factors. On each such slice the extremal problem becomes
a finite optimization over the exponents, and we solve it exactly in
both integer and real variables.

Write \(\omega(n)\) for the number of distinct prime factors of \(n\).
For each \(m\), let \(g_{\integers}(m)\) be the largest log-normalized
divisor sum among integers \(n\geq5041\) with \(\omega(n)=m\). A support
exchange shows that, for \(m\geq6\), every maximizer uses the first \(m\)
primes. Allowing their exponents to be real numbers at least one gives a
relaxation with value \(g_{\reals}(m)\). We study
\[
    \Igap(m)=g_{\reals}(m)-g_{\integers}(m).
\]

\paragraph{Main result.}
If \(p_m\) is the \(m\)th prime, then
\begin{equation}
\label{e-integrality-gap-intro}
    \Igap(m)
    =
    \frac{2\sqrt2+o(1)}
         {\sqrt{p_m}\log p_m}.
\end{equation}
The result is unconditional. The real optimizer is determined by one
scalar equation. Every integer optimizer is a prefix of the exponent
increments ordered by Nicolas's score \(F(p,k)\), with equal scores
taken as one block. This gives a finite
computation with exact score comparisons. Each integer exponent is one of
the two integers adjacent to its relaxed counterpart. The resulting local
rounding sum transfers to a prime integral, where only the second-exponent
range contributes at first order. The two sides of its rounding switch
each contribute \(\sqrt2/(\sqrt{p_m}\log p_m)\).

\paragraph{Robin's inequality.}
The identity
\(-g_{\integers}(m)=\Igap(m)-g_{\reals}(m)\) separates the rounding loss
from the analytic Mertens term. Under the Riemann hypothesis, Nicolas's
integrated explicit formula makes the relaxed value eventually positive
and bounds its normalized coefficient above by
\(4+\gamma-\log(4\pi)\). This is smaller than the gap coefficient
\(2\sqrt2\), so the integer optimum is eventually below Robin's bound. If
the Riemann hypothesis fails, Nicolas's two-sided Mertens excursions
dominate \eqref{e-integrality-gap-intro}. It follows that the hypothesis
is equivalent to Robin's inequality holding on every sufficiently large
fixed-\(\omega\) slice.

\paragraph{Scope.}
We solve the fixed-support relaxation by one scalar equation, characterize
every integer optimizer as a score-ordered prefix, prove one-step
proximity, and evaluate the resulting rounding loss. Combining this
unconditional calculation with Nicolas's integrated explicit formula
gives the fixed-slice criterion. Nicolas's CA score and shell estimates
also place the calculation in the classical theory. In the problem with
\(\omega(n)\leq m\), first-exponent increments become optional and the
same score ordering characterizes the optimizers. A finite verification
settles the slices with \(m\leq5\), where the cutoff \(n\geq5041\)
blocks the support exchange. The integer
\(10080=2^{5}3^{2}5\cdot7\) holds the bounded-support maximum for
\(6\leq m\leq29\), and for \(m\geq30\) the maximizers are score-ordered
prefixes.
